\section{Introduction}

Persistence Mechanics (PM) develops the architecture, dynamics and operational action of persistent systems \cite{kate_meyer_persistence_2026}. Its requirements are organized into Capacity, Map, Protocol, Governor, Toll and Margin. Paper I applies explicit lattice, readability, exchange and selection postulates to the closed vacuum and derives the $E_8$ substrate and its characteristic integers \cite{meyer_informational_2026}. This paper applies that framework to the supplied vacuum substrate as a dynamical test: which update rules, cost functions and continuum responses can maintain the required distinctions? The general action framework is inherited; the task here is its physical realization, not a second derivation of PM.

The operational action is a proposed cost. Its exponential history weights favor lower-cost histories at positive inverse tolerance. A finite penalty suppresses a history without eliminating it; exact concentration requires a limit. Similarly, an exactly unitary lattice update and an approximate continuum equation are different objects. These distinctions determine the scope of the derivation.

\subsection{Mathematical results and physical identifications}
The Lean package checks several linked results. A first-order Hamiltonian has the stipulated relativistic squared dispersion exactly when its generators obey the corresponding Clifford relations. Local phase covariance requires a transporting link. A scalar potential quadratic in $s=|\phi|^2$ has a stable nonzero minimum under explicit no-starvation and no-runaway conditions precisely when its coefficients have the Mexican-hat signs. Finite survival weights admit a concentration limit; detailed balance and an injective symmetrized transfer matrix give a positive two-tick operator and unitary spectral phase evolution.

The new geometric and fermionic calculations make these questions concrete. The $D_4$ nearest-neighbour graph has triangles, whose summed area tensor is isotropic. The Abelian triangle action has a controlled constant-field continuum limit with a calculable normalization. Two $D_4$ sectors support a candidate Wilson/domain-wall mechanism with complementary chiral zero-mode conditions; separate results address walk twins, speed constraints and a detuning mass polynomial.

The new carrier/event calculation follows Paper I's actual minimal glue-class vectors and an
integral $E_8$ root reflection into a counted primitive cost. A specified comparison readout
then yields a trace score and the cosine phase cost, with an explicit primitive-event
multiplicity. The corresponding link expansion now has a controlled finite-angle remainder.
These results begin the common action/response bridge rather than completing its physical
selection. In particular, the three-dimensional cubic walk and the $D_4$ gauge-loop geometry
have not yet been proved to be one microscopic dynamics.

Each theorem has an explicit model boundary. The lattice geometry does not by itself select a physical gauge group, an occupation measure, a full update schedule or a numerical coupling. The two-sector construction is a candidate realization, not a uniqueness theorem derived solely from the sector count. Its exact spectrum, locality properties and response must be connected before numerical constants are inferred.

\subsection{Structure and scope}
\Cref{sec:entropic_action,sec:entropic_lagrangian} state the cost model, survival interpretation and continuum assumptions. \Cref{sec:map_capacity,sec:protocol,sec:governor} present the formalized fermion, gauge and scalar results. \Cref{sec:toll_margin} records the proposed gravitational interpretation and the premises still needed for it. \Cref{sec:synthesis,sec:falsifiability} separate the proved constraints from the next completion tests. The appendices distinguish finite-dimensional quantum reconstruction from a measurement theory and record the uncompleted scattering and magnetic-moment calculations.

The aim is to derive physical laws and their constants from persistence constraints. The current results establish parts of that program, with assumptions exposed. They do not yet establish an exact, parameter-free Standard Model plus General Relativity or the exclusion of all other operators and particle content.
